\chapter{Tu escritorio}\label{cap:06-escritorio}

El \termino{escritorio} es lo que ves después de entrar: el fondo, la barra
con el menú y el reloj, y las ventanas de los programas. Este capítulo
explica XFCE, el que viene marcado de fábrica, y al final resume las
diferencias de los demás.

\section{XFCE de un vistazo}\label{sec:xfce}

\captura{xfce-escritorio}{El escritorio XFCE de Lázaro}

XFCE trae dos barras, que XFCE llama \termino{paneles} \cite{xfce-docs}:

\begin{description}
  \item[El panel de arriba.] A la izquierda, el menú \menu{Aplicaciones}
    con todos los programas, ordenados por tipo (Internet, Oficina,
    Multimedia, Sistema\ldots). En el medio, un botón por cada ventana
    abierta: haz clic para ir a ella o para minimizarla. A la derecha, la
    \termino{bandeja}: los iconos de la red, el volumen, la batería, el
    aviso de actualizaciones y el reloj.
  \item[El panel de abajo.] Un pequeño \enquote{muelle} con accesos a los
    programas más usados: mostrar el escritorio, la terminal, la carpeta
    personal, el navegador y la búsqueda de programas.
\end{description}

\subsection{Abrir un programa}
\begin{pasos}
  \item Haz clic en \menu{Aplicaciones}, arriba a la izquierda.
  \item Pasa el ratón por una categoría, por ejemplo \menu{Oficina}.
  \item Haz clic en el programa.
\end{pasos}

\begin{consejo}
Más rápido: pulsa \keys{Alt + F2} (o la lupa del panel de abajo), escribe
las primeras letras del programa, por ejemplo \texttt{fire} para Firefox, y
pulsa \keys{Enter}.
\end{consejo}

\subsection{Las ventanas}
Cada ventana tiene arriba a la derecha tres botones: \emph{minimizar}
(la esconde en el panel), \emph{maximizar} (ocupa toda la pantalla) y
\emph{cerrar}. Para moverla, arrástrala desde la barra del título. Para
cambiar su tamaño, arrastra un borde o una esquina. Para pasar de una
ventana a otra con el teclado: \keys{Alt + Tab}.

\subsection{Poner un programa en el panel o en el escritorio}
Abre el menú \menu{Aplicaciones}, haz clic derecho sobre el programa y
elige \menu{Añadir al panel} o \menu{Añadir al escritorio}.

\subsection{Cambiar el fondo, la letra y los colores}
\begin{itemize}
  \item \textbf{Fondo:} clic derecho en el escritorio, \menu{Configuración del
        escritorio}. Elige una imagen; para usar una foto tuya, cambia la
        carpeta de arriba.
  \item \textbf{Tamaño de letra y tema claro u oscuro:}
        \menu{Aplicaciones > Configuración > Apariencia}. En la pestaña
        \menu{Estilo} eliges el tema; en \menu{Tipos de letra}, el tamaño.
  \item \textbf{Todo lo demás:} \menu{Aplicaciones > Configuración >
        Gestor de configuración} reúne todos los ajustes en una ventana
        (capítulo~\ref{cap:14-ajustes}).
\end{itemize}

\begin{consejo}
¿Te cuesta leer? Sube el tamaño de letra a 12 o 13 en \menu{Apariencia >
Tipos de letra}, y en Firefox usa \keys{\ctrl + +} para agrandar la página.
\end{consejo}

\subsection{Apagar, reiniciar y cerrar sesión}
En el menú \menu{Aplicaciones}, abajo, está \menu{Salir}. Ofrece
\emph{Apagar}, \emph{Reiniciar}, \emph{Cerrar la sesión} y, en portátiles,
\emph{Suspender}. Apaga siempre así: cortar la corriente de golpe puede
dañar archivos que se estaban guardando.

\section{Los otros escritorios}

Si elegiste otro escritorio al instalar, todo lo de este libro sigue
valiendo: los programas son los mismos. Cambian el aspecto y los nombres
de algunos menús.

\begin{description}
  \item[LXQt] Muy parecido a XFCE pero con una sola barra, abajo, como en
    Windows. El menú está en el botón de la izquierda; la configuración, en
    \menu{Preferencias > Configuración de LXQt} \cite{lxqt-wiki}. Su
    gestor de archivos se llama \emph{PCManFM-Qt}.
  \item[MATE] Dos paneles como XFCE. Arriba, los menús \menu{Aplicaciones},
    \menu{Lugares} y \menu{Sistema}; la configuración está en
    \menu{Sistema > Preferencias} \cite{mate-docs}. Su gestor de archivos se
    llama \emph{Caja}.
  \item[KDE Plasma] Una barra abajo, como Windows, y un menú que se abre con
    el botón de la izquierda o con la tecla \keys{Windows}. Todo se ajusta
    en \emph{Preferencias del sistema}. Los archivos se manejan con
    \emph{Dolphin} \cite{kde-userbase}.
  \item[GNOME] Una barra fina arriba. Pulsa \keys{Windows} (o
    \emph{Actividades}, arriba a la izquierda) para ver las ventanas
    abiertas y buscar programas escribiendo su nombre. La configuración
    está en \emph{Configuración}; los archivos, en \emph{Archivos}
    \cite{gnome-help}.
  \item[Sway] No tiene menús: todo se hace con el teclado. Las ventanas
    se acomodan solas una al lado de la otra. Lo básico: \keys{Windows +
    Enter} abre una terminal, \keys{Windows + D} busca programas,
    \keys{Windows + Mayús + Q} cierra la ventana. Recomendado solo para
    quien ya lo conoce \cite{sway-wiki}.
  \item[XFCE Wayland] Se ve y se usa igual que XFCE. Usa una tecnología
    gráfica más nueva, todavía en prueba: si algo falla, XFCE es la opción
    segura.
\end{description}

\begin{nota}
En los escritorios KDE Plasma y LXQt la tienda de aplicaciones se llama
\emph{Discover}; en los demás, \emph{Software} (capítulo~\ref{cap:13-programas}).
\end{nota}

\fuentes
